% GENERATED FILE. Edit canonical lessons, then run npm run generate:books.
% canonical-source-hash: fnv1a64:9b44821050eed3e7
% canonical-lessons: MR-C344-matbhed, MR-C344-sahmat-hone, MR-C344-aksep-ghene, MR-C344-visangti, MR-C344-drstikon

\chapter{Agreeing and Disagreeing}
\label{ch:mr-agreeing-and-disagreeing}
\hlchaptermodality{344}

\begin{chapteropening}
\textbf{By the end of this chapter:} \emph{I can say that I agree, object or see things differently.}

Complete the last lesson of chapter 344: I can say that I agree, object or see things differently.
\end{chapteropening}

\section[\texorpdfstring{\mr{मतभेद} (matbhed)}{मतभेद (matbhed)}]{\mr{मतभेद} (matbhed) — a disagreement}

\label{lesson:MR-C344-matbhed}



\begin{quote}
Before the new one: say the Marathi for to put forward, then the Marathi for proof.
\end{quote}

\subsection*{You'll want to know: \mr{मतभेद}}
\textbf{\mr{मतभेद}} — \emph{matbhed} — \textquotedblleft{}a disagreement\textquotedblright{}.

\begin{quote}
\mr{भावांमध्ये} \mr{पैशांवरून} \mr{मतभेद} \mr{झाले}.
\end{quote}

\begin{quote}
\emph{bhāvāmmadhye paiśāṁvarūn matbhed jhāle.}
\end{quote}

\begin{quote}
The brothers disagreed over money.
\end{quote}

Masculine.

Mostly plural: \mr{मतभेद} \mr{होणे}, to disagree; \mr{मत} + \mr{भेद}, opinion + split.

\begin{cousinweb}
\textbf{\mr{मतभेद}} (\emph{matbhed}) is built from two tatsamas, Sanskrit \emph{mata}, \textquotedblleft{}thought, opinion\textquotedblright{}, and \emph{bheda}, \textquotedblleft{}splitting, difference\textquotedblright{} (Monier-Williams).
\end{cousinweb}

\begin{grammarlens}[title={What you've built}]
One more word for giving reasons and opinions.
\end{grammarlens}

\subsection*{Your turn}
\begin{itemize}
\raggedright
  \item \emph{Say it:} \emph{matbhed}
  \item \emph{Say it:} \emph{matbhed}, once more
  \item \emph{Say it:} \emph{bhāvāmmadhye paiśāṁvarūn matbhed jhāle.} — The brothers disagreed over money.
  \item \emph{Recall:} say the Marathi for to put forward, then the Marathi for proof, then say \emph{matbhed} again
\end{itemize}

\begin{tcolorbox}[breakable,colback=teal!4,colframe=teal!35!black,title={Before you move on}]
What does \mr{मतभेद} mean? (\textquotedblleft{}A disagreement\textquotedblright{}.) Say it once more.
\end{tcolorbox}
\section[\texorpdfstring{\mr{सहमत} \mr{होणे} (sahmat hoṇe)}{सहमत होणे (sahmat hoṇe)}]{\mr{सहमत} \mr{होणे} (sahmat hoṇe) — to agree}

\label{lesson:MR-C344-sahmat-hone}



\begin{quote}
Before the new one: say the Marathi for proof, then the Marathi for a disagreement.
\end{quote}

\subsection*{You'll want to know: \mr{सहमत} \mr{होणे}}
\textbf{\mr{सहमत} \mr{होणे}} — \emph{sahmat hoṇe} — \textquotedblleft{}to agree\textquotedblright{}.

\begin{quote}
\mr{मी} \mr{तुमच्या} \mr{मताशी} \mr{सहमत} \mr{आहे}.
\end{quote}

\begin{quote}
\emph{mī tumacyā matāśī sahmat āhe.}
\end{quote}

\begin{quote}
I agree with your view.
\end{quote}

Past (he …): \mr{तो} \mr{सहमत} \mr{झाला}, \emph{to sahmat jhālā}.

The person or view agreed with takes -\mr{शी}: \mr{तुमच्याशी} \mr{सहमत}.

\begin{cousinweb}
\textbf{\mr{सहमत} \mr{होणे}} (\emph{sahmat hoṇe}) joins \mr{सहमत}, built from Sanskrit \emph{saha}, \textquotedblleft{}together\textquotedblright{}, and \emph{mata}, \textquotedblleft{}opinion\textquotedblright{} (Monier-Williams), to \mr{होणे}, \textquotedblleft{}to become\textquotedblright{}.
\end{cousinweb}

\begin{grammarlens}[title={What you've built}]
One more word for giving reasons and opinions.
\end{grammarlens}

\subsection*{Your turn}
\begin{itemize}
\raggedright
  \item \emph{Say it:} \emph{sahmat hoṇe}
  \item \emph{Say it:} \emph{sahmat hoṇe}, once more
  \item \emph{Say it:} \emph{mī tumacyā matāśī sahmat āhe.} — I agree with your view.
  \item \emph{Recall:} say the Marathi for proof, then the Marathi for a disagreement, then say \emph{sahmat hoṇe} again
\end{itemize}

\begin{tcolorbox}[breakable,colback=teal!4,colframe=teal!35!black,title={Before you move on}]
What does \mr{सहमत} \mr{होणे} mean? (\textquotedblleft{}To agree\textquotedblright{}.) Say it once more.
\end{tcolorbox}
\section[\texorpdfstring{\mr{आक्षेप} \mr{घेणे} (ākṣep gheṇe)}{आक्षेप घेणे (ākṣep gheṇe)}]{\mr{आक्षेप} \mr{घेणे} (ākṣep gheṇe) — to object}

\label{lesson:MR-C344-aksep-ghene}



\begin{quote}
Before the new one: say the Marathi for a disagreement, then the Marathi for to agree.
\end{quote}

\subsection*{You'll want to know: \mr{आक्षेप} \mr{घेणे}}
\textbf{\mr{आक्षेप} \mr{घेणे}} — \emph{ākṣep gheṇe} — \textquotedblleft{}to object\textquotedblright{}.

\begin{quote}
\mr{अरुणने} \mr{या} \mr{योजनेवर} \mr{आक्षेप} \mr{घेतला}.
\end{quote}

\begin{quote}
\emph{aruṇne yā yojnevar ākṣep ghetlā.}
\end{quote}

\begin{quote}
Arun objected to this plan.
\end{quote}

Past, with -\mr{ने} (he …): \mr{त्याने} \mr{आक्षेप} \mr{घेतला}, \emph{tyāne ākṣep ghetlā}.

What is objected to takes -\mr{वर}: \mr{योजनेवर} \mr{आक्षेप} \mr{घेणे}.

\begin{cousinweb}
\textbf{\mr{आक्षेप} \mr{घेणे}} (\emph{ākṣep gheṇe}) joins the tatsama \mr{आक्षेप}, Sanskrit \emph{ākṣepa}, \textquotedblleft{}throwing at; objection\textquotedblright{} (Monier-Williams), to \mr{घेणे}, \textquotedblleft{}to take\textquotedblright{}.
\end{cousinweb}

\begin{grammarlens}[title={What you've built}]
One more word for giving reasons and opinions.
\end{grammarlens}

\subsection*{Your turn}
\begin{itemize}
\raggedright
  \item \emph{Say it:} \emph{ākṣep gheṇe}
  \item \emph{Say it:} \emph{ākṣep gheṇe}, once more
  \item \emph{Say it:} \emph{aruṇne yā yojnevar ākṣep ghetlā.} — Arun objected to this plan.
  \item \emph{Recall:} say the Marathi for a disagreement, then the Marathi for to agree, then say \emph{ākṣep gheṇe} again
\end{itemize}

\begin{tcolorbox}[breakable,colback=teal!4,colframe=teal!35!black,title={Before you move on}]
What does \mr{आक्षेप} \mr{घेणे} mean? (\textquotedblleft{}To object\textquotedblright{}.) Say it once more.
\end{tcolorbox}
\section[\texorpdfstring{\mr{विसंगती} (visaṅgtī)}{विसंगती (visaṅgtī)}]{\mr{विसंगती} (visaṅgtī) — an inconsistency}

\label{lesson:MR-C344-visangti}



\begin{quote}
Before the new one: say the Marathi for to agree, then the Marathi for to object.
\end{quote}

\subsection*{You'll want to know: \mr{विसंगती}}
\textbf{\mr{विसंगती}} — \emph{visaṅgtī} — \textquotedblleft{}an inconsistency\textquotedblright{}.

\begin{quote}
\mr{त्याच्या} \mr{दोन} \mr{उत्तरांत} \mr{विसंगती} \mr{होती}.
\end{quote}

\begin{quote}
\emph{tyācyā don uttarānt visaṅgtī hotī.}
\end{quote}

\begin{quote}
There was an inconsistency between his two answers.
\end{quote}

Feminine.

Feminine; the adjective is \mr{विसंगत}, inconsistent.

\begin{cousinweb}
\textbf{\mr{विसंगती}} (\emph{visaṅgtī}) is a learned word built from Sanskrit \emph{vi-}, \textquotedblleft{}apart\textquotedblright{}, and \emph{saṁgati}, \textquotedblleft{}coming together\textquotedblright{} (Monier-Williams): things that do not fit together.
\end{cousinweb}

\begin{grammarlens}[title={What you've built}]
One more word for giving reasons and opinions.
\end{grammarlens}

\subsection*{Your turn}
\begin{itemize}
\raggedright
  \item \emph{Say it:} \emph{visaṅgtī}
  \item \emph{Say it:} \emph{visaṅgtī}, once more
  \item \emph{Say it:} \emph{tyācyā don uttarānt visaṅgtī hotī.} — There was an inconsistency between his two answers.
  \item \emph{Recall:} say the Marathi for to agree, then the Marathi for to object, then say \emph{visaṅgtī} again
\end{itemize}

\begin{tcolorbox}[breakable,colback=teal!4,colframe=teal!35!black,title={Before you move on}]
What does \mr{विसंगती} mean? (\textquotedblleft{}An inconsistency\textquotedblright{}.) Say it once more.
\end{tcolorbox}
\section[\texorpdfstring{\mr{दृष्टिकोन} (dṛṣṭikon)}{दृष्टिकोन (dṛṣṭikon)}]{\mr{दृष्टिकोन} (dṛṣṭikon) — a standpoint; an outlook}

\label{lesson:MR-C344-drstikon}



\begin{quote}
Before the new one: say the Marathi for to object, then the Marathi for an inconsistency.
\end{quote}

\subsection*{You'll want to know: \mr{दृष्टिकोन}}
\textbf{\mr{दृष्टिकोन}} — \emph{dṛṣṭikon} — \textquotedblleft{}a standpoint; an outlook\textquotedblright{}.

\begin{quote}
\mr{प्रत्येकाचा} \mr{दृष्टिकोन} \mr{वेगळा} \mr{असतो}.
\end{quote}

\begin{quote}
\emph{pratyekācā dṛṣṭikon vegḷā asto.}
\end{quote}

\begin{quote}
Everyone's outlook is different.
\end{quote}

Masculine.

\mr{दृष्टी} + \mr{कोन}, sight + angle: the angle one looks from.

\begin{cousinweb}
\textbf{\mr{दृष्टिकोन}} (\emph{dṛṣṭikon}) is a learned compound built from Sanskrit \emph{dṛṣṭi}, \textquotedblleft{}sight, view\textquotedblright{}, and \emph{koṇa}, \textquotedblleft{}corner, angle\textquotedblright{} (Monier-Williams).
\end{cousinweb}

\begin{grammarlens}[title={What you've built}]
One more word for giving reasons and opinions.
\end{grammarlens}

\subsection*{Your turn}
\begin{itemize}
\raggedright
  \item \emph{Say it:} \emph{dṛṣṭikon}
  \item \emph{Say it:} \emph{dṛṣṭikon}, once more
  \item \emph{Say it:} \emph{pratyekācā dṛṣṭikon vegḷā asto.} — Everyone's outlook is different.
  \item \emph{Recall:} say the Marathi for to object, then the Marathi for an inconsistency, then say \emph{dṛṣṭikon} again
\end{itemize}

\begin{tcolorbox}[breakable,colback=teal!4,colframe=teal!35!black,title={Before you move on}]
What does \mr{दृष्टिकोन} mean? (\textquotedblleft{}A standpoint; an outlook\textquotedblright{}.) Say it once more.
\end{tcolorbox}
